\chapter{Introduction\sv{Inledning}}
\label{chap:intro}
\begin{instructions}
The full instructions on what to include under the different subheadings will be found on Blackboard. This template is mostly for \textbf{formatting} reference.

\fcolorbox{red}{red!10}{\parbox{0.75\textwidth}{\RaggedRight If there are any discrepancies between the
instructions on Blackboard and this template, the instructions on Blackboard should be followed.}}

The text in the template (like this paragraph) is placed in \texttt{instructions} environments. Write your own text \emph{outside} them. With the class option \texttt{nohints} all instructions and the green Swedish heading names disappear, so you can see your own text and its real length.

If you need to use references you should include them in the file \texttt{thesis.bib} and you should refer to them as \cite{wu2011, overleaf, oetiker2001not}. They appear automatically in the list of references (Bibliography). If they appear separate, it will be with comma between, like these two books~\cite{wu2011, oetiker2001not}. Period should be \emph{after} the reference brackets.

In this chapter you should organise your text into the following sections:
\end{instructions}

\section{Background\sv{Bakgrund}}
\label{sec:bgd}
\begin{instructions}
Give the reader an overview of the context for the thesis.
\end{instructions}

\section{Problem Statement\sv{Problemformulering}}
\label{sec:problem}
\begin{instructions}
Tell the reader who or what has a problem and what is that problem in the context you described.
\end{instructions}

\section{Purpose and goals\sv{Syfte och mål}}
\label{sec:purpose}
\begin{instructions}
Tell the reader what artifact (that you will design and build) can solve the problem.
\end{instructions}

\section{Requirements\sv{Krav}}
\label{sec:req}
\begin{instructions}
Tell the reader about
\begin{itemize}
    \item functional requirements (\textbf{what} should the artifact  do?)
    \item non-functional requirements (\textbf{how} fast or \textbf{how} accurate or \textbf{how} \ldots should the artifact do what it has to do?)
\end{itemize}
Tell the reader about any parts you plan to leave out (if applicable).
\end{instructions}

\section{Related Work\sv{Relaterade arbeten}}
\label{sec:related_work}
\begin{instructions}
Tell the reader the outcome or result of scientifically published work from similar projects or products. You shall write in your own words what they did and cite them properly. Focus on the parts relevant to your project. It does not have to be exactly the same as your project, can be a similar tool but another relevant technique or part of your (or their) project. The reader shall clearly understand (after the problem statement and the purpose) that this forms the knowledge base for your project.
\end{instructions}
